\chapter{Proof Details}\label{app:proofs}

\section{Anchored circle and inverse response}
Let $Z=x+iy\ne0$ and let $Z_0=x_0>0$. The circle with diameter $[0,x_0]$ is
\begin{equation}
 \left(x-\frac{x_0}{2}\right)^2+y^2=\left(\frac{x_0}{2}\right)^2.
 \label{eq:app_circle}
\end{equation}
Expansion gives
\begin{equation}
 x^2+y^2-x_0x=0.
 \label{eq:app_delta_zero}
\end{equation}
Since
\begin{equation}
 \frac1Z=\frac{x-iy}{x^2+y^2},
 \label{eq:app_inverse}
\end{equation}
we have
\begin{equation}
 \RePart\!\left(\frac1Z\right)=\frac{x}{x^2+y^2}.
 \label{eq:app_inverse_real}
\end{equation}
Therefore \cref{eq:app_delta_zero} is equivalent to
\begin{equation}
 1=x_0\RePart\!\left(\frac1Z\right).
 \label{eq:app_classifier}
\end{equation}
At the anchor, $1/Z_0=1/x_0$, completing the proof of \cref{eq:rig_theorem}.

If $1/Z=a+ib$, then $x=a/(a^2+b^2)$ and $|Z|^2=1/(a^2+b^2)$. With $x_0=1/a_0$,
\begin{align}
 \Delta
 &=\frac1{a^2+b^2}
 -\frac1{a_0}\frac{a}{a^2+b^2}\\
 &=\frac{a_0-a}{a_0(a^2+b^2)},
\end{align}
which is \cref{eq:defect_ab}.

\section{Commutant of the shear representation}
For the $120$ degree shear representation matrix
\begin{equation}
 R=\begin{pmatrix}
 -1/2&-\sqrt3/2\\
 \sqrt3/2&-1/2
 \end{pmatrix}
\end{equation}
and a generic real matrix
\begin{equation}
 M=\begin{pmatrix}r&s\\t&u\end{pmatrix},
\end{equation}
the equation $MR-RM=0$ has two independent scalar constraints:
\begin{equation}
 r-u=0,
 \qquad
 s+t=0.
\end{equation}
Thus $M=rI-sJ$ with $J^2=-I$. The independent symbolic audit obtains the same solution with no import from the project response code.

\section{Curvature at the parity anchor}
Write
\begin{equation}
 a(d)=a_0+\frac12a_2d^2+O(d^4),
 \qquad
 b(d)=b_1d+O(d^3),
\end{equation}
and
\begin{equation}
 Z(d)=\frac{a(d)-ib(d)}{a(d)^2+b(d)^2}.
\end{equation}
Parity implies $x'(0)=0$ and $y''(0)=0$. Direct differentiation gives
\begin{equation}
 y'(0)=-\frac{b_1}{a_0^2}
\end{equation}
and
\begin{equation}
 x''(0)=-\frac{a_2}{a_0^2}-\frac{2b_1^2}{a_0^3}.
\end{equation}
Using the signed curvature definition \cref{eq:curvature_general},
\begin{equation}
 \kappa(0)=
 \frac{-y'(0)x''(0)}{|y'(0)|^3}.
\end{equation}
For $b_1>0$, this reduces to
\begin{equation}
 \kappa(0)=-2a_0-\frac{a_0^2a_2}{b_1^2}.
\end{equation}
Changing the orientation changes the overall sign, not the magnitude of the curvature or the dependence on $a_2$.

\section{Quartic drift witness}
For $a(d)=a_0+\epsilon d^4$ and $b(d)=cd$, the second derivative $a''(0)$ vanishes, so the achiral curvature equals the rigid circle value. The exact defect is nevertheless
\begin{equation}
 \Delta(d)=
 -\frac{\epsilon d^4}
 {a_0[(a_0+\epsilon d^4)^2+c^2d^2]}.
\end{equation}
For $a_0>0$, $\epsilon\ne0$, and any nonzero $d$ for which the denominator is finite, the defect is nonzero. This proves that achiral curvature is not a complete global membership test.

\section{Published kinetic circle}
With $x=A/(1+u^2)$ and $y=-Au/(1+u^2)$,
\begin{equation}
 x^2+y^2
 =\frac{A^2+A^2u^2}{(1+u^2)^2}
 =\frac{A^2}{1+u^2}=Ax.
\end{equation}
The kinetic circle follows without approximation after the published coefficients are accepted as the starting expressions.

For the eigenvalue
\begin{equation}
 \lambda=-\frac{D}{2}+i\frac{\sqrt3D}{3}d,
\end{equation}
we have
\begin{align}
 -\frac1{4\lambda^*}
 &=\frac{1}{2D}\frac{1-i(2\sqrt3/3)d}{1+(4/3)d^2}\\
 &=\bar\eta+i\eta_H,
\end{align}
which proves the spectral relation used by the exact test.

\section{OU stationary covariance and Green Kubo integral}
For $K=aI+bJ$ and $C=AI$,
\begin{align}
 KC+CK^T
 &=A(aI+bJ)+A(aI-bJ)\\
 &=2AaI.
\end{align}
If $a>0$, this covariance is positive definite. Since the eigenvalues $a\pm ib$ have positive real part,
\begin{equation}
 \int_0^\infty \e^{-Kt}\,\dd t=K^{-1}.
\end{equation}
Multiplying by $C$ yields \cref{eq:ou_gk}. No numerical quadrature is needed.

\section{Same spectrum, different response matrix element}
Similarity preserves the characteristic polynomial, so the matrices $L_1$ and $L_2$ in \cref{eq:same_spectrum_matrices} have identical eigenvalues. Explicitly,
\begin{equation}
 L_2=\begin{pmatrix}1&1\\0&2\end{pmatrix},
 \qquad
 L_2^{-1}=\begin{pmatrix}1&-1/2\\0&1/2\end{pmatrix}.
\end{equation}
For $g=(1,0)$ and $h=(0,1)^T$,
\begin{equation}
 gL_1^{-1}h=0,
 \qquad
 gL_2^{-1}h=-\frac12.
\end{equation}
Thus an unordered eigenvalue list cannot determine the resolvent matrix element without overlap information.
